% SGA 3, 第 VI_B 篇, 第 1 轮英文重构底本的中译.
% 范围: 命题 8.2、定义 8.2.1 及命题 8.3 全部.
% 权威底本: Polo--Gille Exp6B-13oct24.pdf, 组件第 62--63 页
% (合卷第 403--404 页; 重版页眉第 392--393 页; 原版边栏页码第 376 页).
% 命题 8.4 始于组件第 63 页, 不在本组件内.
% OCR 与 Reinhold 英文 Markdown 仅为对照材料. 本组件无图.

\medskip
\phantomsection\label{prop:sga3-VIB-8.2}%
\noindent\textbf{命题 8.2.---}
\emph{设 \(\mathcal C\) 是一个范畴,
\(T\) 是 \(\mathcal C\) 上的一个拓扑,
\(G\) 是 \(\mathcal C\) 上的一个群层,
\(n\) 是一个满足 \(n\geq0\) 的整数. 以下条件等价:}
\begin{enumerate}
\item[\textup{\SGAInlineRefTarget{prop:sga3-VIB-8.2:item:i}(i)}]\phantomsection\label{concl:sga3-VIB-8.2-i}
  \emph{置 \(K_0=G\), 对 \(p\geq1\), 置}
  \[
    K_p=\operatorname{Comm}(K_{p-1},K_{p-1})
    \qquad
    \bigl(\textup{分别地, }\
      K_p=\operatorname{Comm}(G,K_{p-1})\bigr).
  \]
  \emph{则 \(K_n\) 是单位群预层.}

\item[\textup{\SGAInlineRefTarget{prop:sga3-VIB-8.2:item:ii}(ii)}]\phantomsection\label{concl:sga3-VIB-8.2-ii}
  \emph{置 \(K'_0=G\), 对 \(p\geq1\), 置}
  \[
    K'_p=\operatorname{Comm}_T(K'_{p-1},K'_{p-1})
    \qquad
    \bigl(\textup{分别地, }\
      K'_p=\operatorname{Comm}_T(G,K'_{p-1})\bigr).
  \]
  \emph{则 \(K'_n\) 是单位群层.}

\item[\textup{\SGAInlineRefTarget{prop:sga3-VIB-8.2:item:iii}(iii)}]\phantomsection\label{concl:sga3-VIB-8.2-iii}
  \emph{存在一个由 \(G\) 的正规群子层组成的序列}
  \[
    G=H_0\supset H_1\supset\cdots\supset H_n
  \]
  \emph{使得对每个 \(i\), 商层 \(H_i/H_{i+1}\) 都交换
  \textup{(分别地, 在 \(G/H_{i+1}\) 中是中心的)},
  且 \(H_n\) 是单位层.}
\end{enumerate}

显然 \(K_n\subset K'_n\); 所以
\textup{\SGARefLink{prop:sga3-VIB-8.2:item:ii}{(ii)}} 蕴含
\textup{\SGARefLink{prop:sga3-VIB-8.2:item:i}{(i)}}.
证明 \textup{\SGARefLink{prop:sga3-VIB-8.2:item:i}{(i)}} 蕴含
\textup{\SGARefLink{prop:sga3-VIB-8.2:item:ii}{(ii)}}. 有
\[
  K'_1=\operatorname{Comm}_T(G,G)=K_1^\flat,
\]
并且利用\hyperref[lem:sga3-VIB-8.1.2]{引理 8.1.2} 作归纳, 得到对每个 \(p\), 有
\[
  K'_p=K_p^\flat
\]
所以若 \(K_n\) 是单位预层,
则 \(K'_n=K_n^\flat\) 是单位层.

最后, 条件 \textup{\SGARefLink{prop:sga3-VIB-8.2:item:ii}{(ii)}} 与
\textup{\SGARefLink{prop:sga3-VIB-8.2:item:iii}{(iii)}}
由\hyperref[lem:sga3-VIB-8.1.1]{引理 8.1.1} 可知等价.

\medskip
\phantomsection\label{def:sga3-VIB-8.2.1}%
\noindent\textbf{定义 8.2.1.---}
\emph{这些条件成立时, 称层 \(G\) 是\textup{\(n\) 级可解的}
\textup{(分别地, \(n\) 级幂零的)}.
若存在一个整数 \(n\) 使这些条件成立,
则称 \(G\) 是\textup{可解的}\textup{(分别地, 幂零的)}.}

注意, 由条件 \textup{\SGARefLink{concl:sga3-VIB-8.2-i}{(i)}},
这不依赖于拓扑 \(T\).

\medskip
\phantomsection\label{prop:sga3-VIB-8.3}%
\noindent\textbf{命题 8.3.---}
\emph{设 \(k\) 是域, \(S\) 是非空 \(k\)-概形,
\(\Omega\) 是 \(k\) 的代数闭扩域,
\(G\) 是有限型平滑 \(k\)-群. 以下条件等价:}
\begin{enumerate}
\item[\textup{\SGAInlineRefTarget{prop:sga3-VIB-8.3:item:i}(i)}]\phantomsection\label{concl:sga3-VIB-8.3-i}
  \emph{\(G\) 是 \(n\) 级可解的\textup{(分别地, \(n\) 级幂零的)}.}

\item[\textup{(ii)}]\phantomsection\label{concl:sga3-VIB-8.3-ii}
  \emph{\(G\times_kS\) 是 \(n\) 级可解的\textup{(分别地, \(n\) 级幂零的)}.}

\item[\textup{(iii)}]\phantomsection\label{concl:sga3-VIB-8.3-iii}
  \emph{群 \(G(\Omega)\) 是 \(n\) 级可解的\textup{(分别地, \(n\) 级幂零的)}.}

\item[\textup{(iv)}]\phantomsection\label{concl:sga3-VIB-8.3-iv}
  \emph{置 \(K_0=G\), 对 \(p\geq1\), 考虑 \(k\)-群}
  \[
    K_p=(K_{p-1},K_{p-1})
    \qquad
    \bigl(\textup{分别地, }\ K_p=(G,K_{p-1})\bigr)
  \]
  \emph{\textup{(参见 \hyperref[def:sga3-VIB-7.2.2]{7.2.2})}.
  则 \(K_n\) 是单位 \(k\)-群.}
\end{enumerate}

条件 \textup{\SGARefLink{prop:sga3-VIB-8.3:item:i}{(i)}} 与
\textup{\SGARefLink{concl:sga3-VIB-8.3-i}{(ii)}} 的等价性,
由\hyperref[prop:sga3-VIB-8.2]{命题 8.2} 得到,
因为换位子的群预层的构造与基变换可交换
(IV, \SGARefLink{sga3:iv:proposition:4:1:3}{4.1.3}).

\textup{\SGARefLink{prop:sga3-VIB-8.3:item:i}{(i)}} 与
\textup{\SGARefLink{concl:sga3-VIB-8.3-iii}{(iv)}} 的等价性,
由以下事实得到: 根据 \hyperref[prop:sga3-VIB-7.8]{7.8}, \(k\)-群
\[
  K_p=(K_{p-1},K_{p-1})
  \qquad
  \bigl(\textup{分别地, }\ K_p=(G,K_{p-1})\bigr)
\]
表出层
\[
  \operatorname{Comm}_T(K_{p-1},K_{p-1})
  \qquad
  \bigl(\textup{分别地, }\
    \operatorname{Comm}_T(G,K_{p-1})\bigr),
\]
其中 \(T\) 是 \textup{(fppf)} 或 \textup{(fpqc)} 拓扑.

要证明条件 \textup{\SGARefLink{concl:sga3-VIB-8.3-ii}{(iii)}} 与
\textup{\SGARefLink{concl:sga3-VIB-8.3-iv}{(iv)}} 的等价性,
可以假设 \(k=\Omega\); 此时结果由 \hyperref[cor:sga3-VIB-7.10]{7.10} 得到.

% 中文翻译检查点：本组件全部正文、证明、脚注及源注已译；源139行。
% 机械核对：行内数学52处、陈列数学8处及全部引用键守恒（数学中的自然语言已译）。
% 状态：完成源文覆盖和逐段自校；尚未经独立合卷编译与页面审查。
